\section{Threats to validity and limitations}
\label{sec:limitations}
\paragraph{Evaluation population and dependence}
The 58 eligible targets are not a random deployment sample. Eligibility requires evaluation attacks, sufficient benign calibration data, and enough peers; larger-budget cohorts additionally select for abundant retained training data. Targets can share peer models, captures, collection conditions, and attack-generation procedures. The corpus-stratified bootstrap conditions on the observed corpus groups and does not fully model these dependencies. Three learning seeds assess a fixed design on the same cohort, which was also used during development; they do not sample new corpora. Subgroup and ablation intervals are exploratory, without simultaneous coverage. These boundaries follow established concerns about selection and laboratory evaluation in security machine learning~\cite{arp2024pitfalls}.

\paragraph{Information availability and preprocessing}
Dataset annotations identify the benign onboarding roles. Obtaining equally trustworthy observations in deployment remains an assumption, alongside labeled peer attacks and compatible feature schemas. Additional reference and calibration data and full-training-role pooled standardization make $N$ a local fitting budget rather than a total onboarding budget. Gotham's complete-second aggregates can include attack or unlabeled packets, allow later packets to influence earlier rows, and share raw contributors across role boundaries (Section~\ref{sec:experiment}). This dependency differs from direct use of attack labels in target fitting. Corpus-specific sampling, deduplication, and source ordering further limit chronological interpretation; separate client and label streams do not establish a network-wide prospective cutoff.

\paragraph{Calibration and distribution change}
Calibration controls empirical block rates, without establishing future FPR bounds. The low-alpha allocation exception, reference-tail extrapolation, and final OR without recalibration are part of the implemented rule. Target-level FPR heterogeneity remains substantial under this rule. No online drift detector, adaptive reference update, or recalibration policy was evaluated, and future traffic need not share the calibration distribution.

\paragraph{Comparison and attribution}
The primary comparator changes a complete local branch, including detector composition and calibration. Removal ablations also recalibrate the retained experts. They do not isolate every detector's causal contribution or establish optimality of the allocation. Contemporary external competitors, including MCDE, were not executed under the same protocol; the related-work comparison establishes methodological distinctions, not performance superiority. Attack families can occur in peer development data, so the target information restriction is not uniformly an unseen-attack-family test. Transfer is within compatible study strata, not across corpus schemas.

\paragraph{Security and operational deployment}
The study assumes trustworthy peers and does not test adversarial-peer behavior, poisoning, evasion, encrypted-feature availability under operational constraints, or attacks on the monitor. Model combination and centralized pooled preprocessing establish no privacy guarantee. End-to-end latency, throughput, memory, energy, communications, incident handling, and online deployment were not measured. The shared benchmark pipeline trains the comparator's autoencoder even though PACE excludes its output; algebraic model sizes cannot substitute for measurements of an eventual deployment implementation.
